\scanpage{33}
\mspage{31}

Soit $c$ un \deleted{\illegible{1}} chemin dans $\unsure{M|\mathcal{M}}(X)$, et
\[
  \deleted{\unsure{c\colon}} X\unsure{\mathord{-}|\longrightarrow}\underline{\mathrm{Ch}}.(\deleted{\unsure{X}}M)
\]
le foncteur qu'il définit, \unsure{composons-le} avec les foncteurs source-but \inserted{$\underline{\mathrm{Ch}}(\deleted{\unsure{X}}M)\unsure{\mathord{-}|\longrightarrow}\deleted{\unsure{X}}M\times\deleted{\unsure{X}}M$} on trouve un foncteur \deleted{$X\longrightarrow M$}
\[
 X\longrightarrow M\times M
\]
i.e. un objet \inserted{$(F_0,F_1)$} de $M(X)\times M(X)$, \deleted{d'où \unsure{une|un} \unsure{dé.}} je dis que \deleted{\unsure{ce sont deux|ce sont les}} $F_0$ et $F_1$ sont resp. la source et le but du chemin $c$. De façon un peu plus précise, on a commutativité (stricte) dans le diagramme de foncteurs
\diagram{\begin{tikzcd}[ampersand replacement=\&,column sep=large,row sep=large]
 \underline{\mathrm{Ch}}_{\unsure{\Sigma|Z}_X}(M(X)) \arrow[r,"\sim"] \arrow[d,"{\sigma_{M(X)},\beta_{M(X)}}"] \& \underline{\mathrm{Hom}}^!(X,\underline{\mathrm{Ch}}_\Sigma(M)) \arrow[d,"{\text{déduit de }\sigma_M,\beta_M}"] \\
 M(X)\times M(X) \arrow[r,"\sim"] \& \underline{\mathrm{Hom}}^{\deleted{\unsure{!|?}}}(X,M\times M)
\end{tikzcd}}
où les foncteurs $\sigma$ sont les foncteurs source, les foncteurs $\beta$ \deleted{\unsure{sont}} les foncteurs but.

Si on se donne arbitrairement les objets $F_0,F_1$ dans $M(X)$, donc un foncteur
\[
 X\longrightarrow M\times M
\]
$\underline{\mathrm{Ch}}_{\Sigma_X}(\deleted{\unsure{d}}F_0,F_1)$ s'identifie donc à la \deleted{\unsure{sous}} \inserted{catégorie} des foncteurs \inserted{``admissibles'' \inserted{i.e. se factorisant par un $\underline{\mathrm{Ch}}_{\Sigma,n}$}} $X\xrightarrow{\unsure{c|C}}\underline{\mathrm{Ch}}_\Sigma(M)$ qui \uline{\unsure{relèvent|relient}}
% archive sheet number: 32
